\chapter{38}

\section{Montag, 27. August}



Louis wachte auf. Von draussen erreichte ihn unregelmässiges, helles Glockenschellen. 
Er liess seine Augen langsam durch den Raum wandern, setzte sich auf und schaute aus dem
Fenster. Vor ihm lagen das satte Grün der Wiesenfläche und ein
komplett klarer Himmel. Ihm war, als habe er eine neue
Landschaft vor sich, frisch umgezogen, in feierlicher Kleidung. Erst
jetzt glaubte er, die Schönheit des Ortes wirklich zu sehen.
Ciro durfte hier Musiker sein, auch unfähiger Liebender, mit einer
queren Kindheit und einem abhanden gekommenen Vater, einzig mit Louis und
Giorgia durfte er nichts zu tun haben.
Dennoch beschlich ihn ein unangenehmes Gefühl. Was Manuela und Joh nach
gestern Nacht wohl von ihm hielten? Er war rettungslos in sich
zusammengefallen. Und spätestens jetzt schämte er sich. Die beiden
suchten einen \emph{Zusenn u}nd hatten einen weinenden Kerl bekommen. Es
musste ihm gelingen, sich zu rehabilitieren. Sie sollten überzeugt sein,
in ihm einen taffen Mitarbeiter gewonnen zu haben.
Heute wollte er Maman anrufen.

\sectionbreak

Beim Betreten des Küchenbereichs fiel Louis am Ende der Treppe ein
Zettel auf, der auf dem Boden lag. Handschriftlich und mit grossen,
schwungvollen Buchstaben war eine Mitteilung darauf notiert. Er hob ihn
auf und las.

«Guten Morgen, Ciro, hoffe, du hast prächtig geschlafen. Bin kurz weg.
Gegen 10:00 Uhr zurück. Nimm in der Küche alles, was du brauchst,
Manuela.»

Es lag weder an den Worten noch an Manuelas Handschrift. Es war vielmehr
diese Selbstverständlichkeit, wie die Buchstaben auf diesem Notizzettel
in seiner Hand lagen. Berührt schaute er um sich. Auch draussen war
keiner. Joh musste, wie er erwähnt hatte, in Brig an einer Sitzung sein.
Louis war somit alleine im Haus.
Sie schienen ihm zu vertrauen.
Er wandte sich der Kaffeemaschine zu. Es fühlte sich gut an, sich in
dieser Ruhe einen Cappuccino zu brauen. Einen echten, schaumigen Kaffee.
Mit der Tasse in der Hand lehnte er sich an die Breitseite der
Kochinsel. Er liess seinen Blick langsam durch den Raum gleiten. Neben
einem schwarzen, freistehenden Kühlschrank war eine übergrosse
Magnettafel an der Wand angebracht. Zwischen Unmengen von Fotos und
Ansichtskarten hingen kleine beschriftete Zettel. Interessiert näherte
er sich dem bunten Durcheinander.
Die Fotos zeigten Manuela und Joh, mal zusammen, mal einzeln, junge
Männer, oft in Gruppen, wenige Bilder mit Kindern und Jugendlichen und
an den verschiedensten Orten. Einige der Bilder schienen auf der Alp
aufgenommen zu sein. Da waren sie, diese Schwarznasenschafe. Es machte
den Eindruck, als würden die Tiere ihn direkt ansehen. Die Bilder
wirkten wie Portraits. Es waren fast ausschliesslich Frontaufnahmen von
Schafs-Köpfen.

Eines der grösseren Bilder zog ihn besonders an. Louis ging näher
heran und legte den Kopf schräg. Manuela sass im Gras und schaute direkt
in die Kamera. Auf ihrem Schoss lag anschmiegsam der Kopf eines
wuscheligen Schafes, den sie auf der einen Seite mit der Hand leicht
umschloss. Das Tier hatte die Augen geschlossen, als schliefe es. Hinter
ihr kniete Joh und hielt Manuela mit einem Arm seitlich umfasst. Er
hatte die Augen gesenkt. In seinem andern Arm lag lachend eine junge
Frau. Sie schaute neckisch und direkt zu Joh hinüber. Ihre lag in
Manuelas freier Hand. Jetzt erst entdeckte Louis noch ein weiteres, nur
zur Hälfte sichtbares, männliches Gesicht hinter Joh. Mit
zusammengekniffenen Augen blickte Louis zwischen den beiden jungen
Gesichtern hin und her. Irgendwie glichen sie sich.
Sie sahen alle zufrieden aus. Wie eine moderne Bauernfamilie mit Schaf, dachte er. Dann stutzte er. Eigentlich war es unmöglich, dass sie zusammengehörten. Nach Johs Aussagen waren er und
Manuela doch erst seit zehn Jahre zusammen.

Louis’ Blick wanderte weiter. Die Notizen auf den vereinzelten Zetteln
zwischen Bildern und Ansichtskarten erinnerten ihn an Valentinas
Sammlung. Unterhalb eines Fotos, das Joh zeigte, jünger, verwegen
dreinschauend mit Bart und langem Haar, hing ein Zitat. Ausgedruckt auf
einen schmalen Streifen Papier.

\emph{«Du hast die Wahl. Zu ertrinken oder an Land zu schwimmen. Es
kommt auf deine Entscheidung an.»} Unterhalb stand
handschriftlich und mit kleinen Buchstaben geschrieben \emph{«Ben, dein
Freund, 2008».} Und weiter unten, mit einem andern Stift und in anderer
Schrift notiert, \emph{«vorausgesetzt man kann schwimmen»}. 

Ganz am Rand entdeckte er ein schwarz-weiss Portrait von Manuela, daneben einen hellblau beschrifteten Zettel. 
\emph{«Was der andere an dir ablehnt, ist oft genau das, was ihm an sich fehlt.»}  Louis holte Manuelas Morgennotiz. Er verglich die beiden Schriften. Sie waren identisch. Was sie damit genau sagen wollte? Er würde sie bei Gelegenheit fragen.

Mit der Tasse in der Hand stellte sich Louis in den Türrahmen des
Hauseingangs. Der Tag lag vor ihm wie eine Wundertüte, die darauf wartete, geöffnet zu werden. Wie still es hier war. Kein Hupen, kein geschäftiges Plappern, keine ratternde Strassenbahn, kein
Handyklingeln. Keine Menschenmengen, nichts als diese alles einnehmende
Ruhe. Louis fühlte sich angenehm fremd. Vielleicht musste man sich
ebenso an Stille gewöhnen wie an das emsige Treiben in einer Grossstadt.
Dabei spürte er eine peinigende Sehnsucht nach Mailand, nach seiner Familie,
nach Freunden, seinem Ristorante. Und nach Giorgia.

«Hallo, Ciro.»

Manuelas Ruf klang über den Vorplatz durch die offene Tür zu ihm hoch.
Er hob die Hand zum Gruss. Dass sie ihn Ciro nannte, hörte sich bereits
gewohnt an. Es war verrückt.
Manuela schritt auf das Haus zu. Ihre Wangen waren gerötet. Das leichte
Sommerkleid flatterte um ihre Beine. Sie streckte den Hals und lachte
ihm entgegen. Beladen mit Taschen links und rechts steuerte sie direkt
den Tisch vor dem Haus an. Noch bevor Louis unten war und ihr die beiden
Taschen abnehmen konnte, hievte sie sie mit einem «Uff» beide auf die
Tischplatte und strich sich dann eine lange Haarsträhne aus dem Gesicht.

«Geht schon, danke dir,» ausser Atem rieb sie sich die Hände, «ich habe
ein grandioses Frühstück für uns. Oder hast du bereits gegessen?»

«Nein, ich habe mir erst einen Kaffee gemacht.»

Wieder steckte ihn ihr Strahlen an.

«Ich trag dir die Taschen rein.»

Louis griff sich die beiden prall gefüllten Stoffsäcke und ging vor
Manuela in die Küche hoch. Jetzt erst nahm er die Leere in seinem Magen
wahr. Ihr Frühstücksangebot kam ihm gerade recht. 
Bald sassen sie sich am Tisch draussen gegenüber. Sie bissen herzhaft in
Oliven-Brötchen, die noch ofenwarm in ihren Händen lagen und genossen
sie wie kleine Delikatessen. Einige Nahrungsmittel, erzählte Manuela,
bezogen sie von einem benachbarten Hof. Der lag zu Fuss eine
Viertelstunde nach unten und eine halbe hinauf entfernt. Der
nächstgelegene Hof der \emph{Sunnewaid}.
Louis’ Interesse für die \emph{Sunnewaid} und seine beiden Betreiber war
über die Stunden ungebremst weiter gewachsen.

«Wie lange seid ihr denn schon hier auf der \emph{Sunnewaid}?»

Louis goss sich ein Glas Wasser ein. Manuelas Augen bekamen einen
durchscheinenden Glanz. Es war offenkundig, dass sie ihm gerne antworten
wollte. Sie drehte sich ganz zu ihm um und legte ihre Handflächen auf
den Tisch.

«Wir hatten unglaubliches Glück. Joh hatte einen Onkel, der kinderlos
geblieben und nie verheiratet war. Das war Onkel Theo, zugleich Johs
Pate. Dem gehörten das Haus und das grosse Grundstück. Joh war als Kind
immer mal wieder hier oben. Theo starb vor, Moment, jetzt sind es an die
acht oder neun Jahre.»

Manuela trank einen Schluck Kaffee. Sie räusperte sich geheimnisvoll.

«Als Erwachsener hatte Joh nur mehr losen Kontakt zu Theo, hat er mir
später erzählt. Gemocht hatte er ihn gerne, aber Johs Leben war damals
ziemlich chaotisch. Das müsste er dir selbst erzählen - . Wie auch immer,
kurze Zeit nach Theos Beerdigung, kam ein Brief des Vermögenszentrums
ins Haus geflattert. Joh sollte für die Testamentseröffnung seines
Onkels erscheinen. Ich erinnere mich noch genau, wie perplex er aussah.
Als er zurückkam, war er komplett durch den Wind. Und dann, Ciro, mit
einem Mal veränderte sich Johs Leben. Und meins auch, aber das wurde
erst später klar.»

Louis konnte seinen Blick nicht von Manuelas Gesicht abwenden.

«Theo hatte dem Vermögenszentrum die Verwaltung seines Erbes
anvertraut. Seinem Neffen Johannes schenkte er darin den Hof mit seinen zwei Hektaren
Land. Schon das alleine war eine unerwartete Überraschung. Das Verrückte
daran war, dass Theo keine Bankschuld hatte. Der Hof war urplötzlich in
Johs alleinigem Besitz, zwar renovationsbedürftig, doch abbezahlt,
unfassbar.»

Louis kam vor, als befände sich Manuela mit einem Mal wieder zurück in
dieser Zeit.

«Es kam einem Wunder gleich, als würde dieses Geschenk unsere Beziehung
krönen. Klingt ein bisschen wulstig, das haben wir aber so empfunden,
weil uns dieser Ort, das Haus und die Umgebung ganz neue Möglichkeiten
eröffneten. Joh hatte als Geschiedener und Vater von Zwillingen Alimente
zu bezahlen und wollte seinen Kindern gute gemeinsame Zeiten, Ferien und
so, bieten. Und auf einmal war alles möglich. Mit unserer \emph{Sunnewaid}."

Jetzt sah sie noch ergriffener aus.

«Das war eine der grössten Einschnitte in meinem Leben, etwas, das ich
weder erhofft noch gewünscht hatte, das wäre mir nie eingefallen. Joh
natürlich auch nicht. Wir waren erst ganz kurze Zeit ein Paar. Pochten
auf die Freiheit. Beide hatten wir eine eigene Wohnung. Es ging
schliesslich um mich, weil Joh seine Entscheidung davon abhängig machte,
ob ich miteinstieg. Es war ein Wagnis, partnerschaftlich und auch
finanziell gesehen. Eine grundsätzliche Renovation plus die
Erbschaftssteuer kamen als weitere Herausforderungen dazu. Als Neffe
galt Joh als Verwandter zweiten Grades, nicht ohne, wie hoch die Steuer
war. Heute glaube ich, die wirklich grossen Dinge geschehen, wenn sie
Platz haben, wenn ich bereit bin, ja zu sagen zu etwas Neuem, das sich
mir bietet, zu einer Chance, die ich packen oder ablehnen kann,» sie
atmete aus, «und, Ciro, da habe ich verstanden, was es heisst, eine
bedeutsame Entscheidung zu treffen. Und die brauchte es, von Joh, wie
auch von mir. Er hatte eine sichere Stelle beim Kanton, ich in der
Rechtsmedizin. Ich habe sehr gut verdient, und schliesslich,» sie
runzelte die Stirn und sah Louis nun direkt an, «schliesslich, als die
Visionen wie von selbst kamen und wir uns beide gleichermassen
wohlfühlten darin, gab es kein Halten mehr. Ich hatte einiges zur Seite
gelegt und beteiligte mich finanziell. Mir war klar, es war gewagt.»

Manuelas Augen funkelten.

«Eine der besten Entscheidungen meines Lebens.»

Mit einem Schmunzeln hob sie die Kaffeetasse und trank einen Schluck.

Louis erinnerte sich an Johs Worte zu Haus und Hof vom Vortag. Nach und
nach ergänzten sich seine einzelnen Bilder zu einem Ganzen. Daher also
musste sein Staunen über diesen Ort kommen, der ihm schon beim ersten
Anblick speziell, wenn nicht geradezu, wundersam, vorgekommen war. Die
beiden hatten den Hof bewusst nach ihren Bedürfnissen gestaltet. Sie
lebten offenbar ihre eigenen Träume ebenso wie die gemeinsamen.
Louis stützte unterdessen einhändig sein Kinn auf und staunte Manuela 
stumm an.
Sie rührte gedankenversunken in der Kaffeetasse, die längst leer
geworden war.

«Beeindruckend,» Louis unterbrach sich, schaute Manuela lange an und fuhr schliesslich fort, 
«schöner kann nicht mal eine erfundene Geschichte sein. Man spürt, dass es bei euch stimmt. Es ist
hier alles anders, hm, schon ein bisschen schräg, oder wie soll ich sagen, ungewohnt, für mich, aber gut.»

«Das freut mich, Ciro,» sie lächelte, «bevor dir was \emph{zu} schräg
erscheint, frag nach, unser Ort soll ein freier sein und bleiben. Wir
leben von ehrlicher Kommunikation. Joh und ich sind überzeugt, dass der
direkte Austausch ein wichtiger Schlüssel ist, Probleme zu meistern.»

Sie klang wie Marco. Louis nickte und fühlte plötzlich ein lästiges
Stechen in der Brust.
